\paragraph{Temperature compensation.}
The baseline-subtraction framework and its compensation strategies are
defined by Croxford et al.\ \cite{croxford2010} (BSS/OBS) and Clarke et
al.\ \cite{clarke2010} for long-term temperature cycling of complex
structures. Physics-based \cite{rose2014} and learning-based
compensators followed. Our method is complementary: rather than
compensating temperature away, we exploit the fact that temperature and
damage generate different event signatures.

\paragraph{Guided-wave ML.}
Feature-based ML reached 96.2\% under leave-group-out cross-validation
on the OGW temperature dataset \cite{schnur2022}, with a
$\sim 6\,^\circ$C temperature-extrapolation limit; convolutional
networks on raw signals can exceed compensation outside the baseline
range \cite{mariani2021}; TinyML deployment on microcontrollers has
been demonstrated \cite{henkmann2025tinyml}. We do not compete on peak
accuracy but on the data-rate/Robustness trade-off.

\paragraph{Data reduction.}
Post-hoc SVD/PCA compression of long-term guided-wave matrices is
quantified in \cite{yang2023compression}; long short-term PCA
\cite{yang2022ltpca} targets the same 4.5-year dataset we use. These
are storage-side transforms, not acquisition-side event sampling; we
compare against them at equal data rate (C3).

\paragraph{Event-based sensing.}
Send-on-delta sampling \cite{miskowicz2006} is established in wireless
sensing and neuromorphic vision, and level-crossing ADCs have been
evaluated for ultrasound at the circuit level, but---to our
knowledge---no prior work applies SoD/event-based acquisition to
guided-wave SHM, nor uses event-stream temporal structure to separate
temperature from damage. This is the gap we fill.
